% Bab 2: Use Case
\chapter{Use Case}
\label{bab:use-case}

\section{Diagram Use Case}

Diagram use case berikut memodelkan interaksi antara aktor utama (\textbf{Pengunjung}) dan sistem eksternal (\textbf{Layanan Rute OSRM} serta \textbf{Backend Mangan Yuk}) terhadap 15 use case pada aplikasi \textbf{ManganYuk!}, konsisten dengan \textit{user flow} pada PRD Bagian 2.

\begin{figure}[H]
\centering
\resizebox{\textwidth}{!}{%
\begin{tikzpicture}[
  font=\sffamily\footnotesize,
  uc/.style={ellipse, draw=black!75, fill=blue!6, thick, minimum width=4.6cm, minimum height=0.82cm, align=center, inner sep=2pt},
  ucopt/.style={ellipse, draw=black!75, dashed, fill=teal!6, thick, minimum width=4.6cm, minimum height=0.82cm, align=center, inner sep=2pt},
  actor/.style={align=center},
  assoc/.style={-, thick, draw=black!65},
  ext/.style={->, dashed, thick, draw=black!55}
]

% Use cases - Kolom Kiri (Core Onboarding, Geofence, Eksplorasi)
\node[uc] (uc01) at (0, 11.2) {\textbf{UC-01} Onboarding \& Izin Lokasi};
\node[uc] (uc02) at (0, 9.9)  {\textbf{UC-02} Posisi \& Status Area};
\node[uc] (uc03) at (0, 8.6)  {\textbf{UC-03} Mode Jelajah \& Entry Event};
\node[uc] (uc04) at (0, 7.3)  {\textbf{UC-04} Kartu POI Proximity (10 m)};
\node[uc] (uc05) at (0, 6.0)  {\textbf{UC-05} Catat Kunjungan (Dwell/Exit)};
\node[uc] (uc06) at (0, 4.7)  {\textbf{UC-06} Kartu POI Multimodal};
\node[uc] (uc07) at (0, 3.4)  {\textbf{UC-07} Narasi Eyes-Free};
\node[uc] (uc12) at (0, 2.1)  {\textbf{UC-12} Matikan Mode Jelajah \& Privasi};

% Use cases - Kolom Kanan (Search, Nav, Rec, Trail, Optional)
\node[uc] (uc08) at (5.8, 11.2) {\textbf{UC-08} Cari Tenant \& Tujuan};
\node[uc] (uc09) at (5.8, 9.9)  {\textbf{UC-09} Navigasi Dalam Aplikasi};
\node[uc] (uc10) at (5.8, 8.6)  {\textbf{UC-10} Rekomendasi Berikutnya};
\node[uc] (uc11) at (5.8, 7.3)  {\textbf{UC-11} Smart Culinary Trail};
\node[ucopt] (uc13) at (5.8, 5.4) {\textbf{UC-13} Paspor Kuliner \& Ulasan (F10)};
\node[ucopt] (uc14) at (5.8, 4.1) {\textbf{UC-14} Akun Pengguna Lokal (F11)};
\node[ucopt] (uc15) at (5.8, 2.8) {\textbf{UC-15} Preferensi Halal (F12)};

% System Boundary
\begin{scope}[on background layer]
  \node[draw=black!80, thick, rounded corners=6pt, fill=gray!4,
        fit=(uc01)(uc12)(uc08)(uc15), inner sep=14pt,
        label={[font=\sffamily\bfseries\small, yshift=-14pt]north:Sistem Aplikasi ManganYuk! (Android)}] (sysbox) {};
\end{scope}

% Stick Figure Macro
\newcommand{\stickman}[3]{% name, x, y
  \begin{scope}[shift={(#2,#3)}]
    \draw[thick] (0,0.45) circle (0.18);
    \draw[thick] (0,0.27) -- (0,-0.25);
    \draw[thick] (-0.25,0.10) -- (0.25,0.10);
    \draw[thick] (0,-0.25) -- (-0.20,-0.62);
    \draw[thick] (0,-0.25) -- (0.20,-0.62);
    \node[below=0.70cm, align=center, font=\sffamily\bfseries\footnotesize] (#1) at (0,0) {#1};
  \end{scope}
}

% Aktor Kiri: Pengunjung
\begin{scope}[shift={(-4.6,6.8)}]
  \draw[thick] (0,0.45) circle (0.18);
  \draw[thick] (0,0.27) -- (0,-0.25);
  \draw[thick] (-0.25,0.10) -- (0.25,0.10);
  \draw[thick] (0,-0.25) -- (-0.20,-0.62);
  \draw[thick] (0,-0.25) -- (0.20,-0.62);
  \node[align=center, font=\sffamily\bfseries\footnotesize] at (0,-0.88) {Pengunjung};
  \coordinate (pengunjung) at (0.35, 0);
\end{scope}

% Aktor Kanan 1: Layanan Rute OSRM
\begin{scope}[shift={(10.4,9.9)}]
  \node[draw=black!75, thick, fill=orange!10, rounded corners=3pt, minimum width=2.4cm, minimum height=1.0cm, align=center, font=\sffamily\bfseries\footnotesize] (osrm) {<<system>>\\Layanan Rute\\OSRM};
\end{scope}

% Aktor Kanan 2: Backend Mangan Yuk
\begin{scope}[shift={(10.4,5.4)}]
  \node[draw=black!75, thick, fill=orange!10, rounded corners=3pt, minimum width=2.4cm, minimum height=1.0cm, align=center, font=\sffamily\bfseries\footnotesize] (backend) {<<system>>\\Backend\\Mangan Yuk};
\end{scope}

% Asosiasi Pengunjung -> UC kiri
\foreach \u in {uc01,uc02,uc03,uc04,uc05,uc06,uc07,uc12} {
  \draw[assoc] (pengunjung) -- (\u.west);
}

% Asosiasi Pengunjung -> UC kanan (jalur rapi via sisi atas/bawah)
\foreach \u in {uc08,uc09,uc10,uc11,uc13,uc14,uc15} {
  \draw[assoc] (pengunjung) -- (\u.west);
}

% Asosiasi Aktor Eksternal
\draw[assoc] (uc09.east) -- (osrm.west);
\draw[assoc] (uc13.east) -- (backend.west);

% Legenda singkat di bawah
\node[anchor=north west, font=\sffamily\scriptsize] at ([yshift=-6pt]sysbox.south west)
  {Keterangan: Garis solid = Requirement Utama (F01--F09); Garis putus-putus = Requirement Tambahan/Opsional (F10--F12).};

\end{tikzpicture}%
}
\caption{Diagram Use Case Aplikasi ManganYuk!}
\label{fig:use-case-diagram}
\end{figure}

\section{Daftar Use Case}

\begin{longtable}{|C{0.10}|L{0.25}|L{0.14}|L{0.13}|L{0.38}|}
\caption{Daftar Use Case}\\
\hline
\cellcolor[HTML]{D9EAF7}\textbf{UC ID} & \cellcolor[HTML]{D9EAF7}\textbf{Nama Use Case} & \cellcolor[HTML]{D9EAF7}\textbf{Aktor} & \cellcolor[HTML]{D9EAF7}\textbf{Related Req.} & \cellcolor[HTML]{D9EAF7}\textbf{Deskripsi Singkat} \\ \hline
\endfirsthead
\hline
\cellcolor[HTML]{D9EAF7}\textbf{UC ID} & \cellcolor[HTML]{D9EAF7}\textbf{Nama Use Case} & \cellcolor[HTML]{D9EAF7}\textbf{Aktor} & \cellcolor[HTML]{D9EAF7}\textbf{Related Req.} & \cellcolor[HTML]{D9EAF7}\textbf{Deskripsi Singkat} \\ \hline
\endhead
UC-01 & Melakukan onboarding dan memberi izin lokasi/notifikasi & Pengunjung & F09, F01, F12 & Pengguna membaca penjelasan privasi lokasi, memberi izin lokasi dan notifikasi, lalu memilih preferensi rasa, jenis kuliner, serta opsi halal. \\ \hline
UC-02 & Melihat posisi dan status di dalam/di luar area & Pengunjung & F01 & Aplikasi menampilkan posisi pengguna di peta MapLibre beserta status di dalam/di luar area dan jarak ke Sudirman Street Food. \\ \hline
UC-03 & Mengaktifkan Mode Jelajah dan menerima notifikasi entry & Pengunjung & F02, F09 & Pengguna mengaktifkan Mode Jelajah dan menerima notifikasi sambutan saat memasuki radius 100 m kawasan. \\ \hline
UC-04 & Menerima kartu POI saat mendekati tenant & Pengunjung & F02, F12 & Aplikasi menampilkan kartu POI dan tawaran audio saat pengguna memasuki radius 10 m tenant, tanpa terpicu berulang. \\ \hline
UC-05 & Mencatat kunjungan otomatis (dwell dan exit event) & Pengunjung & F02 & Aplikasi mencatat kunjungan otomatis bila pengguna berada \(\ge\)3 menit di geofence tenant, dan mencatat exit event saat keluar (> radius + 10 m). \\ \hline
UC-06 & Melihat kartu POI multimodal & Pengunjung & F03 & Pengguna membuka kartu POI berisi foto (jika tersedia), cerita budaya, menu \& harga, jam buka, serta audio narasi dan transkrip. \\ \hline
UC-07 & Mendengarkan narasi tanpa melihat layar (eyes-free) & Pengunjung & F04 & Pengguna mendengarkan prompt dan narasi tenant, serta panduan audio, tanpa harus membuka atau melihat layar. \\ \hline
UC-08 & Mencari tenant dan menetapkan tujuan & Pengunjung & F05, F12 & Pengguna mencari tenant berdasarkan nama/menu/kategori (toleran typo), memakai filter (termasuk opsi halal), dan menetapkan tenant sebagai tujuan. \\ \hline
UC-09 & Navigasi menuju kawasan dan tenant & Pengunjung, Layanan Rute OSRM & F06, F01 & Pengguna dipandu di dalam aplikasi menuju tenant (peta MapLibre, rute jalan kaki OSRM atau fallback garis lurus, jarak, panah kompas, audio), termasuk dari luar kawasan. \\ \hline
UC-10 & Menerima rekomendasi tenant berikutnya & Pengunjung & F07, F12 & Setelah keluar dari tenant yang dikunjungi, aplikasi menampilkan 3 pilihan tenant berikutnya beserta alasannya. \\ \hline
UC-11 & Membuat dan menjalankan Smart Culinary Trail & Pengunjung & F08, F02, F12 & Pengguna mengisi durasi, budget, dan rasa; aplikasi menyusun itinerary trail dan memantau progres secara otomatis. \\ \hline
UC-12 & Mematikan Mode Jelajah dan menjaga privasi lokasi & Pengunjung & F09 & Pengguna dapat mematikan Mode Jelajah kapan saja; pelacakan berhenti dan koordinat tidak dikirim ke backend Mangan Yuk. \\ \hline
UC-13 & Mengumpulkan stempel Paspor Kuliner dan memberi ulasan & Pengunjung, Backend Mangan Yuk & F10 & Pengguna memberi ulasan tenant (rating, menu, teks/foto) untuk memperoleh stempel digital dan membuka badge trail tematik. \\ \hline
UC-14 & Mengelola akun pengguna lokal & Pengunjung & F11 & Pengguna masuk atau mendaftar akun lokal, menyimpan sesi di perangkat, mengubah preferensi pada Profil, dan keluar akun. \\ \hline
UC-15 & Mengatur preferensi tenant halal & Pengunjung & F12 & Pengguna mengaktifkan atau menonaktifkan opsi ``Hanya tenant halal'' agar peta, pencarian, trail, rekomendasi, dan trigger menyaring tenant non-halal. \\ \hline
\end{longtable}

\section{Spesifikasi Use Case}

\newcommand{\ucspec}[9]{%
\begin{longtable}{|L{0.26}|L{0.74}|}
\caption{Spesifikasi Use Case #1}\\
\hline
\cellcolor[HTML]{D9EAF7}\textbf{Atribut} & \cellcolor[HTML]{D9EAF7}\textbf{Detail} \\ \hline
\endfirsthead
\hline
\cellcolor[HTML]{D9EAF7}\textbf{Atribut} & \cellcolor[HTML]{D9EAF7}\textbf{Detail} \\ \hline
\endhead
\textbf{Use Case ID} & #1 \\ \hline
\textbf{Nama Use Case} & #2 \\ \hline
\textbf{Aktor} & #3 \\ \hline
\textbf{Related Requirement} & #4 \\ \hline
\textbf{Deskripsi} & #5 \\ \hline
\textbf{Pre-condition} & #6 \\ \hline
\textbf{Trigger} & #7 \\ \hline
\textbf{Alur utama} & #8 \\ \hline
\textbf{Alur alternatif / eksepsi} & #9 \\ \hline
\end{longtable}
}

% UC-01
\ucspec{UC-01}
{Melakukan onboarding dan memberi izin lokasi/notifikasi}
{Pengunjung}
{F09 Location Privacy \& Consent; F01 Location Awareness; F12 Preferensi Halal}
{Pengguna membaca penjelasan privasi lokasi, memberi izin lokasi dan notifikasi, lalu memilih preferensi kuliner dan opsi halal.}
{Aplikasi baru dipasang (data aplikasi bersih); pengguna telah melewati layar login/registrasi (F11); izin lokasi dan notifikasi belum diberikan.}
{Pengguna masuk ke halaman onboarding untuk pertama kali.}
{1. Sistem menampilkan halaman onboarding (jelajah kuliner, cerita budaya, culinary trail).\newline
2. Pengguna menekan tombol Lanjut.\newline
3. Sistem menampilkan penjelasan penggunaan lokasi (alasan, kapan lokasi dibaca, riwayat lokasi tidak dikirim ke backend Mangan Yuk, koordinat hanya dipakai layanan rute OSRM saat navigasi) dan tombol ``Izinkan Lokasi''.\newline
4. Pengguna menekan ``Izinkan Lokasi'', lalu memilih ``Izinkan'' pada dialog sistem.\newline
5. Sistem meminta izin notifikasi; pengguna mengizinkan.\newline
6. Sistem meminta preferensi rasa, jenis kuliner, dan opsi ``Hanya tenant halal'' (dengan opsi Lewati); pengguna memilih lalu menekan Mulai.\newline
7. Sistem menyimpan preferensi secara lokal dan menampilkan Beranda.}
{4a. Pengguna menekan ``Tolak'' pada dialog izin lokasi: sistem menampilkan penjelasan bahwa trigger otomatis dan navigasi tidak tersedia, dengan tombol Buka Pengaturan dan Lanjut tanpa Lokasi.\newline
4b. Pengguna menekan ``Lanjut tanpa Lokasi'': sistem menampilkan peta tanpa posisi pengguna dan daftar tenant; kartu POI dan audio dapat dibuka manual.\newline
6a. Pengguna menekan Lewati pada preferensi: sistem melanjutkan ke Beranda.\newline\textbf{Post-condition:} Preferensi tersimpan secara lokal; Beranda tampil (mode dengan lokasi atau mode tanpa lokasi).}

% UC-02
\ucspec{UC-02}
{Melihat posisi dan status di dalam/di luar area}
{Pengunjung}
{F01 Location Awareness}
{Aplikasi menampilkan posisi pengguna di peta MapLibre beserta status di dalam/di luar area dan jarak ke Sudirman Street Food.}
{Aplikasi terbuka di Beranda; izin lokasi diberikan; GPS aktif.}
{Pengguna membuka Beranda atau berpindah posisi.}
{1. Sistem membaca posisi pengguna (pembaruan tiap \(\pm\)3 detik atau perpindahan \(\ge\)5 m).\newline
2. Sistem menampilkan titik biru dengan lingkaran akurasi pada peta MapLibre (tile OpenStreetMap) dan denah kawasan.\newline
3. Sistem menampilkan status ``Di dalam/Di luar area Sudirman Street Food''.\newline
4. Sistem memperbarui posisi dan status saat pengguna berjalan.}
{3a. Pengguna berada di luar kawasan: sistem menampilkan status ``Di luar area'' dan jarak ke Sudirman Street Food; pengguna tetap dapat memilih tenant untuk memulai navigasi rute OSRM dari luar kawasan (lihat UC-09).\newline\textbf{Post-condition:} Posisi dan status area tampil sesuai kondisi aktual pengguna.}

% UC-03
\ucspec{UC-03}
{Mengaktifkan Mode Jelajah dan menerima notifikasi entry}
{Pengunjung}
{F02 Location-Based Trigger (Geofence); F09 Location Privacy \& Consent}
{Pengguna mengaktifkan Mode Jelajah dan menerima notifikasi sambutan saat memasuki radius 100 m kawasan.}
{Onboarding selesai; izin lokasi dan notifikasi diberikan; pengguna menuju kawasan.}
{Pengguna mengaktifkan Mode Jelajah dan memasuki radius 100 m kawasan.}
{1. Pengguna mengaktifkan Mode Jelajah.\newline
2. Sistem mengaktifkan pelacakan lokasi di background (\textit{foreground service}) dan menampilkan notifikasi permanen ``Mangan Yuk! sedang memandu''.\newline
3. Pengguna memasuki radius 100 m kawasan.\newline
4. Sistem memeriksa akurasi pembacaan lokasi (\(\le\)30 m).\newline
5. Sistem menampilkan notifikasi entry dan audio sambutan (``Selamat datang di Sudirman Street Food!'').\newline
6. Sistem menawarkan pilihan Rencanakan Trail atau Jelajah Bebas.}
{4a. Akurasi >30 m: sistem mengabaikan pembacaan, menampilkan banner ``Sinyal GPS lemah, trigger otomatis ditunda''; pengguna dapat membuka kartu POI manual dari daftar ``Tenant di sekitar''; trigger otomatis berlanjut setelah akurasi \(\le\)30 m.\newline
6a. Pengguna memilih Jelajah Bebas: peta menampilkan seluruh marker tenant sesuai preferensi halal tanpa tujuan tertentu dan proximity event tetap aktif.\newline\textbf{Post-condition:} Mode Jelajah aktif; kedatangan di kawasan terdeteksi.}

% UC-04
\ucspec{UC-04}
{Menerima kartu POI saat mendekati tenant}
{Pengunjung}
{F02 Location-Based Trigger (Geofence); F12 Preferensi Halal}
{Aplikasi menampilkan kartu POI dan tawaran audio saat pengguna memasuki radius 10 m tenant, tanpa terpicu berulang.}
{Mode Jelajah aktif; izin lokasi diberikan; data 31 POI termuat dari \texttt{pois.json}.}
{Posisi pengguna masuk ke radius proximity tenant (10 m).}
{1. Sistem membaca posisi pengguna dari pembaruan lokasi background.\newline
2. Sistem menghitung jarak ke setiap POI (haversine) dan mengenali POI terdekat.\newline
3. Sistem memeriksa pengguna berada dalam radius 10 m, akurasi \(\le\)30 m, preferensi halal terpenuhi (jika F12 aktif), dan POI belum terpicu/tidak dalam cooldown.\newline
4. Sistem menampilkan kartu POI, getar singkat, dan tawaran audio (``Kamu mendekati [tenant]. Mau dengar ceritanya?'').\newline
5. Pengguna membuka detail tenant atau mendengarkan narasi.}
{3a. POI sudah pernah terpicu dan pengguna masih di area yang sama, atau POI dalam cooldown 30 menit: sistem tidak menampilkan notifikasi ulang.\newline
3b. Akurasi lokasi >30 m: pembacaan diabaikan dan trigger ditunda (lihat UC-03).\newline
3c. Opsi ``Hanya tenant halal'' aktif dan POI berstatus non-halal/belum diverifikasi: proximity event tidak dipicu.\newline\textbf{Post-condition:} POI tercatat sebagai sudah terpicu (cooldown 30 menit); kartu POI tampil.}

% UC-05
\ucspec{UC-05}
{Mencatat kunjungan otomatis (dwell dan exit event)}
{Pengunjung}
{F02 Location-Based Trigger (Geofence)}
{Aplikasi mencatat kunjungan otomatis bila pengguna berada \(\ge\)3 menit di geofence tenant dan mencatat exit event saat keluar.}
{Mode Jelajah aktif; pengguna berada di dalam geofence tenant (radius 10 m).}
{Pengguna berada di dalam geofence tenant.}
{1. Sistem menghitung durasi pengguna di dalam geofence tenant.\newline
2. Setelah \(\ge\)3 menit, sistem mencatat kunjungan (\textit{dwell event}) di perangkat dan menandai tenant sudah dikunjungi.\newline
3. Bila trail aktif, sistem menandai tenant pada trail sebagai selesai.\newline
4. Saat pengguna keluar dari geofence (jarak > radius + 10 m, yakni >20 m), sistem mencatat \textit{exit event}.}
{2a. Pengguna keluar dari geofence sebelum 3 menit (hanya lewat): kunjungan tidak tercatat.\newline\textbf{Post-condition:} Kunjungan tercatat pada riwayat lokal perangkat.}

% UC-06
\ucspec{UC-06}
{Melihat kartu POI multimodal}
{Pengunjung}
{F03 Multimodal Exploration}
{Pengguna membuka kartu POI berisi foto (jika tersedia), cerita budaya, menu \& harga, jam buka, serta audio narasi dan transkrip.}
{Aplikasi terbuka; data 31 POI dibundel di dalam APK.}
{Pengguna memilih tenant dari peta/daftar atau kartu POI muncul dari trigger.}
{1. Pengguna memilih tenant dari peta/daftar atau menerima kartu POI dari trigger.\newline
2. Sistem membuka kartu POI (foto makanan/tenant jika tersedia, kategori dan profil rasa, menu \& harga, jam buka).\newline
3. Pengguna memutar audio narasi TTS (45--90 detik) yang memuat cerita kuliner/budaya; audio dapat dijeda dan dilanjutkan.\newline
4. Pengguna membuka transkrip audio.}
{3a. Tidak ada koneksi internet: seluruh konten dataset, foto lokal, dan audio tetap dapat diakses karena dibundel di APK.\newline\textbf{Post-condition:} Informasi tenant ditampilkan secara lengkap.}

% UC-07
\ucspec{UC-07}
{Mendengarkan narasi tanpa melihat layar (eyes-free)}
{Pengunjung}
{F04 Eyes-Free Interaction}
{Pengguna mendengarkan prompt dan narasi tenant, serta panduan audio, tanpa harus membuka atau melihat layar.}
{Mode Jelajah aktif; izin lokasi dan notifikasi diberikan; layar terkunci atau HP di saku.}
{Proximity event pada tenant tujuan atau tenant yang dilewati.}
{1. Sistem memutar prompt suara (``Kamu mendekati [tenant]. Mau dengar ceritanya?'') dan HP bergetar.\newline
2. Notifikasi menampilkan tombol aksi Dengar/Lewati, termasuk pada lock screen.\newline
3. Pengguna menekan Dengar.\newline
4. Sistem memutar audio narasi tanpa membuka aplikasi.\newline
5. Pada Mode Jalan, narasi diputar otomatis; navigasi memberi panduan audio tiap pengurangan jarak 20 m dan saat tiba.}
{3a. Pengguna menekan Lewati: prompt ditutup, tenant masuk daftar ``Terlewati'', dan tidak dipicu ulang selama cooldown 30 menit.\newline\textbf{Post-condition:} Narasi terdengar (atau dilewati) tanpa pengguna membuka aplikasi.}

% UC-08
\ucspec{UC-08}
{Mencari tenant dan menetapkan tujuan}
{Pengunjung}
{F05 Search \& Destination; F12 Preferensi Halal}
{Pengguna mencari tenant berdasarkan nama/menu/kategori (toleran typo), memakai filter, dan menetapkan tenant sebagai tujuan.}
{Aplikasi terbuka; data 31 POI tersedia secara offline.}
{Pengguna menekan kolom pencarian.}
{1. Pengguna menekan kolom pencarian dan mengetik kata kunci.\newline
2. Sistem menampilkan hasil pencarian toleran typo, diurutkan berdasarkan jarak.\newline
3. Pengguna menerapkan filter (rentang harga, rasa, jenis, status buka).\newline
4. Sistem memperbarui daftar sesuai kombinasi filter (dan menyembunyikan tenant non-halal jika F12 aktif).\newline
5. Pengguna memilih tenant dan menekan tombol Jadikan Tujuan.\newline
6. Sistem masuk ke mode navigasi menuju tenant yang dipilih.}
{2a. Tidak ada tenant yang sesuai dengan kata kunci/filter: sistem menampilkan daftar kosong/pesan tidak ditemukan tanpa crash.\newline\textbf{Post-condition:} Tenant tujuan ditetapkan dan mode navigasi aktif.}

% UC-09
\ucspec{UC-09}
{Navigasi menuju kawasan dan tenant}
{Pengunjung, Layanan Rute OSRM}
{F06 Navigation; F01 Location Awareness}
{Pengguna dipandu di dalam aplikasi menuju tenant tujuan (peta MapLibre, rute jalan kaki OSRM atau fallback garis lurus, jarak, panah kompas, audio), termasuk saat memulai dari luar kawasan.}
{Izin lokasi diberikan; tujuan ditetapkan.}
{Pengguna menekan Jadikan Tujuan atau Mulai Trail.}
{1. Pengguna menetapkan tenant tujuan (dari dalam maupun luar kawasan).\newline
2. Sistem menampilkan rute jalan kaki OSRM di peta MapLibre dari posisi pengguna ke tenant, jarak dalam meter, dan panah arah berbasis kompas.\newline
3. Jika pengguna keluar >25 m dari rute, sistem menghitung ulang rute OSRM.\newline
4. Sistem memberi panduan audio setiap jarak berkurang 20 m.\newline
5. Saat jarak \(\le\)10 m, sistem memutar ``Kamu sudah sampai!'' dan membuka kartu POI tenant tujuan.}
{2a. Layanan OSRM gagal atau tidak ada koneksi internet: sistem menampilkan garis lurus ke tenant tujuan dan navigasi tetap berjalan.\newline
1a. Izin lokasi ditolak: navigasi tidak aktif (mode tanpa lokasi).\newline\textbf{Post-condition:} Pengguna tiba di tenant tujuan dan kartu POI terbuka.}

% UC-10
\ucspec{UC-10}
{Menerima rekomendasi tenant berikutnya}
{Pengunjung}
{F07 Contextual Recommendation; F12 Preferensi Halal}
{Setelah keluar dari tenant yang dikunjungi, aplikasi menampilkan 3 pilihan tenant berikutnya beserta alasannya.}
{Mode Jelajah aktif; minimal satu tenant sudah dikunjungi (dwell tercatat); preferensi tersimpan.}
{Exit event dari geofence tenant yang sudah dikunjungi (jarak >20 m).}
{1. Pengguna berjalan keluar dari geofence tenant yang sudah dikunjungi.\newline
2. Sistem menghitung skor \textit{rule-based} (jarak, status buka, kecocokan rasa, sisa budget trail, preferensi halal, status kunjungan, aturan pasangan rasa).\newline
3. Sistem menampilkan ``3 pilihan berikutnya'' beserta alasan masing-masing.\newline
4. Pengguna memilih tenant dan masuk ke mode navigasi.}
{2a. Tenant tutup, sudah dikunjungi, atau non-halal saat F12 aktif: tenant dikeluarkan dari kandidat rekomendasi.\newline\textbf{Post-condition:} Rekomendasi tampil; pengguna dapat memilih tujuan berikutnya.}

% UC-11
\ucspec{UC-11}
{Membuat dan menjalankan Smart Culinary Trail}
{Pengunjung}
{F08 Smart Culinary Trail Planner; F02 Location-Based Trigger (Geofence); F12 Preferensi Halal}
{Pengguna mengisi durasi, budget, dan rasa; aplikasi menyusun itinerary trail dan memantau progres secara otomatis.}
{Mode Jelajah aktif; pengguna berada di dalam kawasan; preferensi tersimpan.}
{Pengguna memilih Rencanakan Trail.}
{1. Pengguna memilih Rencanakan Trail (dari notifikasi entry atau Beranda).\newline
2. Sistem menampilkan form trail: durasi (30/60/90 menit), budget, preferensi rasa (terisi otomatis dari onboarding).\newline
3. Pengguna memilih durasi 60 menit, budget Rp50.000, lalu menekan Buat Trail.\newline
4. Sistem menampilkan itinerary tenant berurutan (\textit{nearest neighbor}, dengan status halal jika F12 aktif) beserta estimasi total biaya, estimasi waktu, dan rute di peta.\newline
5. Pengguna menekan Mulai Trail; sistem masuk mode navigasi tenant pertama.\newline
6. Setiap dwell event menandai tenant selesai; sistem mengirim pengingat sisa waktu bila dwell melebihi alokasi + 5 menit; di akhir ditampilkan ringkasan trail.}
{6a. Tenant tutup atau pengguna menekan Lewati Tenant: sistem menandai tenant dilewati dan menyusun ulang sisa trail dari posisi pengguna sesuai sisa waktu dan sisa budget (tenant yang sudah selesai/dilewati tidak dipilih lagi).\newline\textbf{Post-condition:} Progres trail tersimpan; ringkasan trail (jumlah tenant, estimasi belanja, durasi, jarak) tampil.}

% UC-12
\ucspec{UC-12}
{Mematikan Mode Jelajah dan menjaga privasi lokasi}
{Pengunjung}
{F09 Location Privacy \& Consent}
{Pengguna dapat mematikan Mode Jelajah kapan saja; pelacakan berhenti dan koordinat tidak dikirim ke backend Mangan Yuk.}
{Mode Jelajah aktif.}
{Pengguna menekan Selesai Jelajah.}
{1. Pengguna menekan tombol Selesai Jelajah.\newline
2. Sistem menghentikan pelacakan lokasi di background.\newline
3. Sistem menghapus notifikasi permanen.\newline
4. Sistem kembali ke Beranda.}
{Selama aplikasi digunakan, tidak ada request ke backend Mangan Yuk yang membawa koordinat pengguna; koordinat hanya dikirim ke layanan rute OSRM saat layar navigasi aktif.\newline\textbf{Post-condition:} Pelacakan background berhenti dan tidak ada koordinat pengguna yang dikirim ke backend Mangan Yuk.}

% UC-13
\ucspec{UC-13}
{Mengumpulkan stempel Paspor Kuliner dan memberi ulasan}
{Pengunjung, Backend Mangan Yuk}
{F10 Paspor Kuliner \& Ulasan Tenant (Gamification)}
{Pengguna memberi ulasan tenant untuk memperoleh stempel digital pada trail tematik dan membuka badge.}
{Pengguna sudah masuk akun lokal (F11); halaman Paspor Kuliner atau detail tenant terbuka.}
{Pengguna menekan Tambah Ulasan.}
{1. Pengguna membuka Paspor Kuliner dari Profil atau menekan Tambah Ulasan pada detail tenant.\newline
2. Pengguna memilih tenant, memberi rating 1--5, memilih menu yang dimakan (milik tenant tersebut), serta mengisi teks ulasan dan foto (opsional).\newline
3. Sistem menyimpan ulasan, mencatat stempel digital untuk tenant tersebut, dan memperbarui progres trail tematik (``Jejak Kuliner Nusantara'' atau ``Manis-manis Bandung'').\newline
4. Sistem menyinkronkan data ke backend Mangan Yuk (\texttt{/v1/passport} dan \texttt{/v1/reviews}).\newline
5. Saat seluruh 5 tenant pada satu trail berstempel, sistem membuka badge trail tersebut.}
{3a. Pengguna memberi ulasan kedua pada tenant yang sama: ulasan baru tersimpan tanpa menambah stempel ganda.\newline
4a. Backend tidak dapat dihubungi: ulasan dan stempel tetap tersimpan di perangkat (AsyncStorage).\newline\textbf{Post-condition:} Ulasan, stempel, dan progres trail tersimpan dan bertahan setelah aplikasi ditutup ulang.}

% UC-14
\ucspec{UC-14}
{Mengelola akun pengguna lokal}
{Pengunjung}
{F11 Akun Pengguna Lokal}
{Pengguna mendaftar atau masuk ke akun lokal, menyimpan sesi secara lokal di perangkat, dan mengelola profil.}
{Aplikasi terpasang di perangkat.}
{Pengguna membuka aplikasi tanpa sesi aktif atau membuka halaman Profil.}
{1. Sistem mengarahkan pengguna yang belum masuk ke halaman login.\newline
2. Pengguna mendaftar atau masuk dengan email, nama pengguna, dan kata sandi (opsional mengaktifkan ``Ingat Saya'').\newline
3. Sistem menyimpan sesi di perangkat melalui SecureStore (tanpa mengirim ke server) dan mengarahkan ke onboarding atau Beranda.\newline
4. Pada halaman Profil, pengguna dapat melihat data akun, ringkasan paspor, mengubah preferensi, atau menekan Keluar.}
{4a. Pengguna menekan Keluar: sistem menghapus sesi dari perangkat dan mengarahkan ke halaman login.\newline\textbf{Post-condition:} Sesi akun lokal tersimpan atau terhapus sesuai aksi pengguna.}

% UC-15
\ucspec{UC-15}
{Mengatur preferensi tenant halal}
{Pengunjung}
{F12 Preferensi Halal}
{Pengguna mengaktifkan atau menonaktifkan opsi ``Hanya tenant halal'' dari onboarding atau Profil.}
{Pengguna berada di halaman onboarding atau Profil.}
{Pengguna mengubah toggle ``Hanya tenant halal''.}
{1. Pengguna mengaktifkan opsi ``Hanya tenant halal''.\newline
2. Sistem menyaring marker dan carousel peta, hasil pencarian (F05), pilihan Smart Culinary Trail (F08), rekomendasi (F07), serta proximity event (F02) agar hanya memproses tenant berstatus halal.\newline
3. Pengguna dapat menonaktifkan kembali opsi tersebut dari halaman Profil untuk menampilkan seluruh tenant.}
{1a. Tenant berstatus non-halal atau belum diverifikasi: disembunyikan dari peta, pencarian, trail, dan rekomendasi serta tidak memicu notifikasi proximity saat opsi aktif.\newline\textbf{Post-condition:} Seluruh fitur terkait mematuhi status preferensi halal pengguna.}
